\documentclass[10pt]{beamer}
\usetheme{metropolis}
\usepackage{graphicx}\usepackage{listings}\usepackage{booktabs}\usepackage{xcolor}
\definecolor{codebg}{HTML}{F5F5F5}\definecolor{codegreen}{HTML}{2E7D32}\definecolor{codegray}{HTML}{757575}\definecolor{codeblue}{HTML}{1565C0}
\lstdefinestyle{rstyle}{language=R,backgroundcolor=\color{codebg},basicstyle=\ttfamily\scriptsize,keywordstyle=\color{codeblue}\bfseries,commentstyle=\color{codegray}\itshape,breaklines=true,frame=single,rulecolor=\color{codegray!30},showstringspaces=false,morekeywords={library,ggplot,aes,geom_line,geom_area,geom_step,scale_x_date,date_breaks,date_labels,theme_minimal,labs,stl,decompose,ts,as.Date,ymd,year,month}}
\lstdefinestyle{pythonstyle}{language=Python,backgroundcolor=\color{codebg},basicstyle=\ttfamily\scriptsize,keywordstyle=\color{codeblue}\bfseries,commentstyle=\color{codegray}\itshape,breaklines=true,frame=single,rulecolor=\color{codegray!30},showstringspaces=false,morekeywords={import,from,as,pd,np,plt,to_datetime,resample,DateFormatter,seasonal_decompose,statsmodels}}
\title{Time as a Visual Dimension}
\subtitle{Workshop 6 --- Module 1}
\author{Prof.Asoc.\ Endri Raco}
\institute{Data Visualization for Data Scientists \\ UNYT Tirana}
\date{2025/26}
\begin{document}

\begin{frame}
  \titlepage
\end{frame}

\begin{frame}{Workshop 6 Overview}
  \begin{center}
    \includegraphics[width=0.78\textwidth]{figures_w06m01/w06_overview}
  \end{center}
\end{frame}

\begin{frame}{Module Roadmap}
  \begin{enumerate}
    \item What makes time special as a visual dimension
    \item Temporal granularity: daily, weekly, monthly, yearly
    \item Time series decomposition: trend + seasonal + residual
    \item The temporal chart taxonomy: five core types
    \item Time-axis pitfalls: spacing, truncation, missing dates
  \end{enumerate}
  \begin{exampleblock}{Learning Outcome}
    You will identify the appropriate temporal granularity for a dataset,
    decompose a time series into trend, seasonal, and residual components,
    and choose the right temporal chart type based on the data structure.
  \end{exampleblock}
\end{frame}

\section{Why Time Is Special}

\begin{frame}{Time as the Organising Axis}
  Time is unique among variables:
  \begin{itemize}
    \item It is \textbf{ordered} and \textbf{continuous} (unlike categories)
    \item It has \textbf{natural units} (second, hour, day, week, month, year)
    \item It enables \textbf{causal reasoning}: before $\rightarrow$ after
    \item It has \textbf{hierarchical granularity}: zoom in/out
  \end{itemize}
  \vspace{6pt}
  Because of these properties, temporal data demands specialised chart types
  (line, area, step) and specialised axis formatting (date breaks, labels,
  proper spacing). A scatter plot of date vs value ``works'' but wastes
  time's natural structure.
  \begin{alertblock}{Pro-Tip}
    The x-axis of a time series chart is \textbf{not a categorical axis}.
    It must use a true date/time scale so that gaps between observations
    reflect real elapsed time. Never convert dates to strings and plot
    as categories.
  \end{alertblock}
\end{frame}

\section{Temporal Granularity}

\begin{frame}{Same Data, Four Granularities}
  \begin{center}
    \includegraphics[width=0.88\textwidth]{figures_w06m01/granularity}
  \end{center}
\end{frame}

\begin{frame}{Choosing Granularity}
  \centering
  \begin{tabular}{llp{4.5cm}}
    \toprule
    \textbf{Granularity} & \textbf{Best For} & \textbf{Risk} \\
    \midrule
    Daily & Short spans ($<$3 months), event detection & Noisy for longer spans \\
    Weekly & Medium spans (3--12 months), smooths weekday effects & Loses within-week patterns \\
    Monthly & Annual comparisons, seasonal analysis & Loses event timing \\
    Quarterly & Financial reporting, macro trends & Very coarse \\
    Yearly & Multi-decade trends, cross-country comparison & Loses all sub-annual patterns \\
    \bottomrule
  \end{tabular}
  \vspace{6pt}
  {\small Rule: start at the \textbf{finest available granularity}, then
  aggregate up. You can always smooth daily $\rightarrow$ monthly, but
  you cannot recover daily detail from monthly data.}
\end{frame}

\begin{frame}[fragile]{Resampling in Code}
  \begin{columns}[T]
    \begin{column}{0.48\textwidth}
      \textbf{R}
\begin{lstlisting}[style=rstyle]
library(lubridate); library(dplyr)

daily <- read_csv("data.csv") |>
  mutate(date = ymd(date))

# Weekly mean
weekly <- daily |>
  mutate(week = floor_date(date,
    "week")) |>
  group_by(week) |>
  summarise(value = mean(value))

# Monthly
monthly <- daily |>
  mutate(month = floor_date(date,
    "month")) |>
  group_by(month) |>
  summarise(value = mean(value))
\end{lstlisting}
    \end{column}
    \begin{column}{0.48\textwidth}
      \textbf{Python}
\begin{lstlisting}[style=pythonstyle]
df["date"] = pd.to_datetime(df["date"])
df = df.set_index("date")

# Weekly mean
weekly = df.resample("W").mean()

# Monthly
monthly = df.resample("ME").mean()

# Quarterly
quarterly = df.resample("QE").mean()

# Custom: 2-week rolling average
rolling_14 = df.rolling("14D").mean()
\end{lstlisting}
    \end{column}
  \end{columns}
\end{frame}

\section{Time Series Decomposition}

\begin{frame}{Trend + Seasonal + Residual}
  \begin{center}
    \includegraphics[width=0.82\textwidth]{figures_w06m01/decomposition}
  \end{center}
\end{frame}

\begin{frame}{Two Decomposition Models}
  \centering
  \begin{tabular}{llp{4.5cm}}
    \toprule
    \textbf{Model} & \textbf{Formula} & \textbf{When to Use} \\
    \midrule
    Additive & $Y = T + S + R$ & Seasonal amplitude is constant (temperature) \\
    Multiplicative & $Y = T \times S \times R$ & Seasonal amplitude grows with trend (retail sales) \\
    \bottomrule
  \end{tabular}
  \vspace{8pt}
  \begin{alertblock}{Pro-Tip}
    If the ``peaks and valleys'' get taller as the trend rises, use
    multiplicative. If they stay the same height regardless of the
    level, use additive. Quick test: plot the data --- if the envelope
    widens over time, it is multiplicative.
  \end{alertblock}
\end{frame}

\begin{frame}[fragile]{Decomposition in Code}
  \begin{columns}[T]
    \begin{column}{0.48\textwidth}
      \textbf{R}
\begin{lstlisting}[style=rstyle]
# Base R
ts_data <- ts(daily$value,
  frequency = 365)
decomp <- decompose(ts_data,
  type = "additive")
plot(decomp)

# STL (robust, flexible)
stl_fit <- stl(ts_data,
  s.window = "periodic")
plot(stl_fit)

# Extract components
trend    <- stl_fit$time.series[,"trend"]
seasonal <- stl_fit$time.series[,"seasonal"]
residual <- stl_fit$time.series[,"remainder"]
\end{lstlisting}
    \end{column}
    \begin{column}{0.48\textwidth}
      \textbf{Python}
\begin{lstlisting}[style=pythonstyle]
from statsmodels.tsa.seasonal import (
    seasonal_decompose, STL)

# Classical decomposition
result = seasonal_decompose(
    df["value"],
    model="additive",
    period=365)
result.plot()

# STL (robust)
stl = STL(df["value"], period=365)
res = stl.fit()
res.plot()

# Extract
trend    = res.trend
seasonal = res.seasonal
residual = res.resid
\end{lstlisting}
    \end{column}
  \end{columns}
\end{frame}

\section{Temporal Chart Taxonomy}

\begin{frame}{Five Core Temporal Chart Types}
  \begin{center}
    \includegraphics[width=0.92\textwidth]{figures_w06m01/taxonomy}
  \end{center}
\end{frame}

\begin{frame}{Choosing the Right Chart}
  \centering
  \begin{tabular}{llll}
    \toprule
    \textbf{Question} & \textbf{Chart} & \textbf{Why} \\
    \midrule
    How does Y change over time? & Line & Continuous trajectory \\
    What is the cumulative total? & Stacked area & Parts sum to whole \\
    When did Y change discretely? & Step & Rate changes, policies \\
    What are the period totals? & Bar (temporal) & Monthly/yearly sums \\
    What is the daily pattern? & Heatmap calendar & Day-level grid \\
    How do many series compare? & Sparklines & Compact, Tufte-style \\
    \bottomrule
  \end{tabular}
\end{frame}

\section{Pitfalls}

\begin{frame}{Time-Axis Pitfalls}
  \begin{center}
    \includegraphics[width=0.88\textwidth]{figures_w06m01/pitfalls}
  \end{center}
  \begin{alertblock}{Pro-Tip}
    Three more pitfalls: (1) \textbf{Truncated y-axis} --- starting at
    a high value exaggerates small changes. (2) \textbf{Missing dates}
    --- gaps in the data should appear as gaps in the line, not as
    connected straight segments. (3) \textbf{Ignoring time zones} ---
    always record and document the time zone.
  \end{alertblock}
\end{frame}

\section{Synthesis}

\begin{frame}{Key Takeaways}
  \begin{enumerate}
    \item Time is \textbf{ordered, continuous, hierarchical, and causal}
          --- demanding specialised chart types and axis formatting.
    \item \textbf{Granularity}: start fine (daily), aggregate up.
          Daily for events, weekly for medium spans, monthly for
          seasonal analysis, yearly for long trends.
    \item \textbf{Decomposition}: $Y = T + S + R$ (additive) or
          $Y = T \times S \times R$ (multiplicative). Use STL for
          robust decomposition. R: \texttt{stl()}. Python:
          \texttt{statsmodels.tsa.seasonal.STL()}.
    \item Five temporal chart types: \textbf{line} (trend), \textbf{area}
          (cumulative), \textbf{step} (discrete changes), \textbf{bar}
          (period totals), \textbf{heatmap calendar} (daily patterns).
    \item \textbf{Never} use categorical spacing for dates. Always
          use a true date/time axis.
  \end{enumerate}
\end{frame}

\begin{frame}{Essential Readings}
  \begin{itemize}
    \item Cleveland, R.~B. et al. (1990). ``STL: A Seasonal-Trend
          Decomposition Procedure Based on Loess.'' \emph{Journal
          of Official Statistics}, 6(1).
    \item Hyndman, R.~J. \& Athanasopoulos, G. (2021). \emph{Forecasting:
          Principles and Practice}, 3rd ed. OTexts.
          \texttt{https://otexts.com/fpp3/}
    \item Wilke, C.~O. (2019). \emph{Fundamentals of Data Visualization},
          Chapter 13 (Visualizing Time Series).
  \end{itemize}
  \vspace{6pt}
  \textbf{Next Module}: Line Charts --- Theory \& Best Practices (W06-M02)
\end{frame}

\end{document}
